%\subsection{聖靈}

\subsection{預備}
\subsubsection{預備总结}
\begin{table}[H]
\centering
\begin{tabular}{l|l|l|l}\hline
\multicolumn{2}{c|}{神的預備} & \multicolumn{2}{c}{人預備的目的是迎接人子,進羔羊婚娶} \\\hline\hline
王為他兒子預備娶親的筵席& 太 22:2-14&預備人子來了& 太 24:32-44\\\hline
王為他兒子預備娶親的筵席& 路 14:15-24&預備油在器皿裡迎接新郎來進去坐席& 太 25:1-13\\\hline
預備地方& 約 14:1-3&預備永不壞的錢囊、用不盡的財寶& 路 12:33-34\\\hline
預備一座城& 來 11:16&腰束帶，燈要點,像僕人等主人從婚筵回來& 路 12:35-40\\\hline
末世要顯現的救恩& 彼前1:5&按时分粮,顺主人的意思行& 路 12:42-48\\\hline
&&預備心中盼望的緣由& 彼前 3:13-16\\\hline
&&新婦自己預備羔羊婚娶& 啟 19:7\\\hline\hline
\end{tabular}
\end{table}
 

\subsubsection{預備:迎接人子來/油/永不壞的錢囊,用不盡的財寶/腰束帶燈要點/進羔羊婚娶}
\begin{itemize}
\item \textbf{預備人子來了} - 馬太福音 24:32-44\\
你們可以從無花果樹學個比方。當樹枝發嫩長葉的時候，你們就知道夏天近了。 33 這樣，你們看見這一切的事，也該知道人子近了，正在門口了。 34 我實在告訴你們：這世代還沒有過去，這些事都要成就。 35 天地要廢去，我的話卻不能廢去。 36 但那日子、那時辰，沒有人知道，連天上的使者也不知道，子也不知道，唯獨父知道。 37 挪亞的日子怎樣，人子降臨也要怎樣。 38 當洪水以前的日子，人照常吃喝嫁娶，直到挪亞進方舟的那日， 39 不知不覺洪水來了，把他們全都沖去。人子降臨也要這樣。 40 那時，兩個人在田裡，取去一個，撇下一個； 41 兩個女人推磨，取去一個，撇下一個。 42 所以，你們要警醒，因為不知道你們的主是哪一天來到。 43 家主若知道幾更天有賊來，就必警醒，不容人挖透房屋，這是你們所知道的。 44 所以，你們也要\textcolor{red}{預備}，因為你們想不到的時候，人子就來了。
\item \textbf{預備油在器皿裡迎接新郎來進去坐席} - 馬太福音 25:1-13\\
那時，天國好比十個童女拿著燈出去迎接新郎。 2 其中有五個是愚拙的，五個是聰明的。 3 愚拙的拿著燈，卻\textcolor{red}{預備}油； 4 聰明的拿著燈，又\textcolor{red}{預備}油在器皿裡。 5 新郎遲延的時候，她們都打盹睡著了。 6 半夜有人喊著說：『新郎來了，你們出來迎接他！』 7 那些童女就都起來收拾燈。 8 愚拙的對聰明的說：『請分點油給我們，因為我們的燈要滅了。』 9 聰明的回答說：『恐怕不夠你我用的，不如你們自己到賣油的那裡去買吧。』 10 她們去買的時候，新郎到了，那\textcolor{red}{預備}好了的同他進去坐席，門就關了。 11 其餘的童女隨後也來了，說：『主啊，主啊，給我們開門！』 12 他卻回答說：『我實在告訴你們：我不認識你們。』 13 所以，你們要警醒，因為那日子、那時辰，你們不知道。
\item \textbf{預備永不壞的錢囊、用不盡的財寶} - 路加福音 12:33-34\\
你們要變賣所有的賙濟人，為自己\textcolor{red}{預備}永不壞的錢囊、用不盡的財寶在天上，就是賊不能近、蟲不能蛀的地方。 34 因為你們的財寶在哪裡，你們的心也在哪裡。

\item \textbf{腰裡束上帶，燈也要點著,好像僕人等候主人從婚姻的筵席上回來} - 路加福音 12:35-40\\
35 「你們腰裡要束上帶，燈也要點著， 36 自己好像僕人等候主人從婚姻的筵席上回來。他來到叩門，就立刻給他開門。 37 主人來了，看見僕人警醒，那僕人就有福了。我實在告訴你們：主人必叫他們坐席，自己束上帶進前伺候他們。 38 或是二更天來，或是三更天來，看見僕人這樣，那僕人就有福了。 39 家主若知道賊什麼時候來，就必警醒，不容賊挖透房屋，這是你們所知道的。 40 你們也要\textcolor{red}{預備}，因為你們想不到的時候，人子就來了。

\item \textbf{按时分粮顺主人的意思行} - 路加福音 12:42-48\\
谁是那忠心有见识的管家，主人派他管理家里的人，按时分粮给他们呢？ 43 主人来到，看见仆人这样行，那仆人就有福了。 44 我实在告诉你们：主人要派他管理一切所有的。 45 那仆人若心里说‘我的主人必来得迟’，就动手打仆人和使女，并且吃喝醉酒， 46 在他想不到的日子，不知道的时辰，那仆人的主人要来，重重地处治他，定他和不忠心的人同罪。 47 仆人知道主人的意思，却不\textcolor{red}{預備}，又不顺他的意思行，那仆人必多受责打。 48 唯有那不知道的，做了当受责打的事，必少受责打。因为多给谁，就向谁多取；多托谁，就向谁多要。

\item \textbf{預備心中盼望的緣由} - 彼得前書 3:13-16\\
你們若是熱心行善，有誰害你們呢？ 14 你們就是為義受苦，也是有福的。不要怕人的威嚇，也不要驚慌， 15 只要心裡尊主基督為聖。有人問你們心中盼望的緣由，就要常做\textcolor{red}{準備}，以溫柔、敬畏的心回答各人， 16 存著無虧的良心，叫你們在何事上被毀謗，就在何事上可以叫那誣賴你們在基督裡有好品行的人自覺羞愧。

\item \textbf{新婦自己預備羔羊婚娶} - 啟示錄 19:7\\
我們要歡喜快樂，將榮耀歸給他，因為羔羊婚娶的時候到了，新婦也自己\textcolor{red}{預備}好了，
\end{itemize}

\subsubsection{神的預備:末世要顯現的救恩/娶親的筵席}

\begin{itemize}
\item 王為他兒子預備娶親的筵席 - 馬太福音 22:2-14\\
天國好比一個王為他兒子擺設娶親的筵席， 3 就打發僕人去，請那些被召的人來赴席，他們卻不肯來。 4 王又打發別的僕人，說：『你們告訴那被召的人：我的筵席已經\textcolor{red}{預備}好了，牛和肥畜已經宰了，各樣都\textcolor{red}{齊備}，請你們來赴席。』 5 那些人不理就走了：一個到自己田裡去，一個做買賣去， 6 其餘的拿住僕人，凌辱他們，把他們殺了。 7 王就大怒，發兵除滅那些凶手，燒毀他們的城。 8 於是對僕人說：『喜筵已經\textcolor{red}{齊備}，只是所召的人不配。 9 所以你們要往岔路口上去，凡遇見的，都召來赴席。』 10 那些僕人就出去，到大路上，凡遇見的，不論善惡都召聚了來，筵席上就坐滿了客。 11 王進來觀看賓客，見那裡有一個沒有穿禮服的， 12 就對他說：『朋友，你到這裡來怎麼不穿禮服呢？』那人無言可答。 13 於是王對使喚的人說：『捆起他的手腳來，把他丟在外邊的黑暗裡，在那裡必要哀哭切齒了。』 14 因為被召的人多，選上的人少。
\item 王為他兒子預備娶親的筵席 - 路加福音 14:15-24\\
同席的有一人聽見這話，就對耶穌說：「在神國裡吃飯的有福了！」 16 耶穌對他說：「有一人擺設大筵席，請了許多客。 17 到了坐席的時候，打發僕人去對所請的人說：『請來吧，樣樣都\textcolor{red}{齊備}了。』 18 眾人一口同音地推辭。頭一個說：『我買了一塊地，必須去看看。請你准我辭了。』 19 又有一個說：『我買了五對牛，要去試一試。請你准我辭了。』 20 又有一個說：『我才娶了妻，所以不能去。』 21 那僕人回來，把這事都告訴了主人。家主就動怒，對僕人說：『快出去到城裡大街小巷，領那貧窮的、殘廢的、瞎眼的、瘸腿的來。』 22 僕人說：『主啊，你所吩咐的已經辦了，還有空座。』 23 主人對僕人說：『你出去到路上和籬笆那裡，勉強人進來，坐滿我的屋子。 24 我告訴你們：先前所請的人，沒有一個得嘗我的筵席！』」
\item 預備地方 - 約翰福音 14:1-3\\
你們心裡不要憂愁，你們信神，也當信我。 2 在我父的家裡有許多住處，若是沒有，我就早已告訴你們了，我去原是為你們\textcolor{red}{預備}地方去。 3 我若去為你們\textcolor{red}{預備}了地方，就必再來接你們到我那裡去，我在哪裡，叫你們也在哪裡。

\item 預備了一座城 - 希伯來書 11:16\\
他們卻羨慕一個更美的家鄉，就是在天上的。所以神被稱為他們的神，並不以為恥，因為他已經給他們\textcolor{red}{預備}了一座城。

\item 末世要顯現的救恩- 彼得前书1:5\\
你們這因信蒙神能力保守的人，必能得著所\textcolor{red}{預備}、到末世要顯現的救恩。

\end{itemize}


\subsection{警醒}
\subsubsection{警醒禱告/做工/看门/等主/記念圣经教导/在真道上立穩/穿戴救恩/不被吞吃/堅固弱者/看守衣服}
\begin{table}[H]
\centering
\begin{tabular}{l|l|l|l}\hline
\multicolumn{2}{c|}{四福音} & \multicolumn{2}{c}{其他} \\\hline\hline
主哪一天來到& 太 24:32-44&警醒記念保罗三年晝夜不住流淚的勸誡& 徒 20:25-31\\\hline
迎接新郎,進去坐席& 太 25:1-13&警醒在真道上立穩，做大丈夫，要剛強&林前16:13 \\\hline
警醒禱告，免得入了迷惑& 太 26:37-41&恆切禱告警醒、感恩&西 4:2-4 \\\hline
警醒禱告/做工/看门，迎接主來& 可 13:33-37&不要睡覺要警醒謹守,穿戴信/愛/得救的盼望&贴前 5:6 \\\hline
警醒禱告，免得入了迷惑& 可 14:32-38&謹慎自守警醒禱告& 彼前4:7\\\hline
束腰點燈,等候主人從婚筵回來& 路 12:35-40&警醒被仇敵吞吃&彼前5:8 \\\hline
常常祈求逃避要来的,站立在人子前& 路 21:34-36&堅固那剩下將要衰微的，& 啟 3:2-3\\\hline
& &警醒看守衣服免得赤身而行&啟 16:15 \\\hline\hline
\end{tabular}
\end{table}

\subsubsection{详细经文}
馬太福音 24:32-44\\
你們可以從無花果樹學個比方。當樹枝發嫩長葉的時候，你們就知道夏天近了。 33 這樣，你們看見這一切的事，也該知道人子近了，正在門口了。 34 我實在告訴你們：這世代還沒有過去，這些事都要成就。 35 天地要廢去，我的話卻不能廢去。 36 但那日子、那時辰，沒有人知道，連天上的使者也不知道，子也不知道，唯獨父知道。 37 挪亞的日子怎樣，人子降臨也要怎樣。 38 當洪水以前的日子，人照常吃喝嫁娶，直到挪亞進方舟的那日， 39 不知不覺洪水來了，把他們全都沖去。人子降臨也要這樣。 40 那時，兩個人在田裡，取去一個，撇下一個； 41 兩個女人推磨，取去一個，撇下一個。 42 所以，你們要\textcolor{red}{警醒}，因為不知道你們的主是哪一天來到。 43 家主若知道幾更天有賊來，就必\textcolor{red}{警醒}，不容人挖透房屋，這是你們所知道的。 44 所以，你們也要預備，因為你們想不到的時候，人子就來了。


馬太福音 25:1-13\\
那時，天國好比十個童女拿著燈出去迎接新郎。 2 其中有五個是愚拙的，五個是聰明的。 3 愚拙的拿著燈，卻預備油； 4 聰明的拿著燈，又預備油在器皿裡。 5 新郎遲延的時候，她們都打盹睡著了。 6 半夜有人喊著說：『新郎來了，你們出來迎接他！』 7 那些童女就都起來收拾燈。 8 愚拙的對聰明的說：『請分點油給我們，因為我們的燈要滅了。』 9 聰明的回答說：『恐怕不夠你我用的，不如你們自己到賣油的那裡去買吧。』 10 她們去買的時候，新郎到了，那預備好了的同他進去坐席，門就關了。 11 其餘的童女隨後也來了，說：『主啊，主啊，給我們開門！』 12 他卻回答說：『我實在告訴你們：我不認識你們。』 13 所以，你們要\textcolor{red}{警醒}，因為那日子、那時辰，你們不知道。

馬太福音 26:36-41\\
你們坐在這裡，等我到那邊去禱告。」 37 於是帶著彼得和西庇太的兩個兒子同去，就憂愁起來，極其難過， 38 便對他們說：「我心裡甚是憂傷，幾乎要死。你們在這裡等候，和我一同\textcolor{red}{警醒}。」 39 他就稍往前走，俯伏在地，禱告說：「我父啊！倘若可行，求你叫這杯離開我！然而，不要照我的意思，只要照你的意思。」 40 來到門徒那裡，見他們睡著了，就對彼得說：「怎麼樣，你們不能同我\textcolor{red}{警醒}片時嗎？ 41 總要\textcolor{red}{警醒}禱告，免得入了迷惑。你們心靈固然願意，肉體卻軟弱了。」


马可福音 13:33-37\\
你们要谨慎，\textcolor{red}{警醒}祈祷，因为你们不晓得那日期几时来到。 34 这事正如一个人离开本家，寄居外邦，把权柄交给仆人，分派各人当做的工，又吩咐看门的\textcolor{red}{警醒}。 35 所以，你们要\textcolor{red}{警醒}，因为你们不知道家主什么时候来，或晚上，或半夜，或鸡叫，或早晨。 36 恐怕他忽然来到，看见你们睡着了。 37 我对你们所说的话，也是对众人说：要\textcolor{red}{警醒}！

马可福音 14:32-38\\
耶穌對門徒說：「你們坐在這裡，等我禱告。」 33 於是帶著彼得、雅各、約翰同去，就驚恐起來，極其難過， 34 對他們說：「我心裡甚是憂傷，幾乎要死。你們在這裡等候，\textcolor{red}{警醒}。」 35 他就稍往前走，俯伏在地，禱告說倘若可行，便叫那時候過去。 36 他說：「阿爸，父啊！在你凡事都能，求你將這杯撤去！然而，不要從我的意思，只要從你的意思。」 37 耶穌回來，見他們睡著了，就對彼得說：「西門，你睡覺嗎？不能\textcolor{red}{警醒}片時嗎？ 38 總要\textcolor{red}{警醒}禱告，免得入了迷惑。你們心靈固然願意，肉體卻軟弱了。」

路加福音 12:35-40\\
35 「你們腰裡要束上帶，燈也要點著， 36 自己好像僕人等候主人從婚姻的筵席上回來。他來到叩門，就立刻給他開門。 37 主人來了，看見僕人\textcolor{red}{警醒}，那僕人就有福了。我實在告訴你們：主人必叫他們坐席，自己束上帶進前伺候他們。 38 或是二更天來，或是三更天來，看見僕人這樣，那僕人就有福了。 39 家主若知道賊什麼時候來，就必\textcolor{red}{警醒}，不容賊挖透房屋，這是你們所知道的。 40 你們也要預備，因為你們想不到的時候，人子就來了。

路加福音 21:34-36\\
你们要谨慎，恐怕因贪食、醉酒并今生的思虑累住你们的心，那日子就如同网罗忽然临到你们， 35 因为那日子要这样临到全地上一切居住的人。 36 你们要时时\textcolor{red}{警醒}，常常祈求，使你们能逃避这一切要来的事，得以站立在人子面前。

使徒行傳 20:25-31\\
我素常在你們中間來往，傳講神國的道，如今我曉得，你們以後都不得再見我的面了。 26 所以我今日向你們證明，你們中間無論何人死亡，罪不在我身上， 27 因為神的旨意我並沒有一樣避諱不傳給你們的。 28 聖靈立你們做全群的監督，你們就當為自己謹慎，也為全群謹慎，牧養神的教會，就是他用自己血所買來的 d。 29 我知道我去之後，必有凶暴的豺狼進入你們中間，不愛惜羊群。 30 就是你們中間也必有人起來，說悖謬的話，要引誘門徒跟從他們。 31 所以你們應當\textcolor{red}{警醒}，記念我三年之久晝夜不住地流淚勸誡你們各人。

哥林多前書 16:13\\
你們務要\textcolor{red}{警醒}，在真道上站立得穩，要做大丈夫，要剛強。

歌羅西書 4:2-4
你們要恆切禱告，在此\textcolor{red}{警醒}、感恩。 3 也要為我們禱告，求神給我們開傳道的門，能以講基督的奧祕——我為此被捆鎖—— 4 叫我按著所該說的話將這奧祕發明出來。

帖撒羅尼迦前書 5:6-11\\
所以，我們不要睡覺，像別人一樣，總要\textcolor{red}{警醒}謹守。 7 因為睡了的人是在夜間睡，醉了的人是在夜間醉。 8 但我們既然屬乎白晝，就應當謹守，把信和愛當做護心鏡遮胸，把得救的盼望當做頭盔戴上。 9 因為神不是預定我們受刑，乃是預定我們藉著我們主耶穌基督得救。 10 他替我們死，叫我們無論\textcolor{red}{醒著}、睡著，都與他同活。 11 所以，你們該彼此勸慰，互相建立，正如你們素常所行的。

彼得前書 4:7\\
萬物的結局近了，所以你們要謹慎自守，\textcolor{red}{警醒}禱告。

警醒仇敵吞吃 - 彼得前書 5:8\\
務要謹守、\textcolor{red}{警醒}，因為你們的仇敵魔鬼如同吼叫的獅子，遍地遊行，尋找可吞吃的人。

啟示錄 3:2,3\\
你要\textcolor{red}{警醒}，堅固那剩下將要衰微的，因我見你的行為在我神面前，沒有一樣是完全的。…

警醒來看守衣服 - 啟示錄 16:15\\
看哪，我來像賊一樣。那\textcolor{red}{警醒}、看守衣服，免得赤身而行，叫人見他羞恥的有福了！



\subsection{谨慎}
\subsubsection{谨慎}
%馬太福音 18:10\\
%這世界有禍了，因為將人絆倒；絆倒人的事是免不了的，但那絆倒人的有禍了！ 8 倘若你一隻手或是一隻腳叫你跌倒，就砍下來丟掉！你缺一隻手或是一隻腳進入永生，強如有兩手兩腳被丟在永火裡。 9 倘若你一隻眼叫你跌倒，就把它剜出來丟掉！你只有一隻眼進入永生，強如有兩隻眼被丟在地獄的火裡。 10 你們要\textcolor{red}{小心}，不可輕看這小子裡的一個。我告訴你們，他們的使者在天上，常見我天父的面。


馬太福音 24:7-14\\
耶穌回答說：「你們要\textcolor{red}{谨慎}，免得有人迷惑你們。 5 因為將來有好些人冒我的名來，說：『我是基督』，並且要迷惑許多人。 6 你們也要聽見打仗和打仗的風聲，總不要驚慌，因為這些事是必須有的，\textcolor{blue}{只是末期還沒有到}。
民要攻打民，國要攻打國，多處必有饑荒、地震， 8 這都是災難 a的起頭。 9 那時，人要把你們陷在患難裡，也要殺害你們；你們又要為我的名被萬民恨惡。10 「那時，必有許多人跌倒，也要彼此陷害、彼此恨惡， 11 且有好些假先知起來，迷惑多人。 12 只因不法的事增多，許多人的愛心才漸漸冷淡了。 13 唯有忍耐到底的，必然得救。 14 這天國的福音要傳遍天下，對萬民作見證，\textcolor{blue}{然後末期才來到}。

马可福音 4:24\\
你們所聽的要\textcolor{red}{留心}。你們用什麼量器量給人，也必用什麼量器量給你們，並且要多給你們。 25 因為有的，還要給他；沒有的，連他所有的也要奪去。


马可福音 8:14-15\\
門徒忘了帶餅，在船上除了一個餅，沒有別的食物。 15 耶穌囑咐他們說：「你們要\textcolor{red}{谨慎}，防備法利賽人的酵和希律的酵。」

马可福音 12:38-40\\
耶穌在教訓之間說：「你們要\textcolor{red}{防備}文士。他們好穿長衣遊行，喜愛人在街市上問他們的安， 39 又喜愛會堂裡的高位、筵席上的首座。 40 他們侵吞寡婦的家產，假意作很長的禱告。這些人要受更重的刑罰！」





马可福音 13:5-13\\
耶稣说：“你们要\textcolor{red}{谨慎}，免得有人迷惑你们。 6 将来有好些人冒我的名来，说‘我是基督’，并且要迷惑许多人。 7 你们听见打仗和打仗的风声，不要惊慌，这些事是必须有的，只是末期还没有到。 8 民要攻打民，国要攻打国，多处必有地震、饥荒，这都是灾难的起头。但你们要\textcolor{red}{谨慎}，因为人要把你们交给公会，并且你们在会堂里要受鞭打，又为我的缘故站在诸侯与君王面前，对他们作见证。 10 然而，福音必须先传给万民。 11 人把你们拉去交官的时候，不要预先思虑说什么，到那时候，赐给你们什么话，你们就说什么；因为说话的不是你们，乃是圣灵。 12 弟兄要把弟兄，父亲要把儿子，送到死地；儿女要起来与父母为敌，害死他们。 13 并且你们要为我的名被众人恨恶，唯有忍耐到底的，必然得救。

马可福音 13:21-23\\
那時，若有人對你們說『看哪，基督在這裡』，或說『基督在那裡』，你們不要信。 22 因為假基督、假先知將要起來，顯神蹟奇事，倘若能行，就把選民迷惑了。 23 你們要\textcolor{red}{謹慎}！看哪，凡事我都預先告訴你們了！


马可福音 13:33-36\\
你们要\textcolor{red}{谨慎}，警醒祈祷，因为你们不晓得那日期几时来到。 34 这事正如一个人离开本家，寄居外邦，把权柄交给仆人，分派各人当做的工，又吩咐看门的警醒。 35 所以，你们要警醒，因为你们不知道家主什么时候来，或晚上，或半夜，或鸡叫，或早晨。 36 恐怕他忽然来到，看见你们睡着了。 37 我对你们所说的话，也是对众人说：要警醒！

路加福音 8:16-18\\
沒有人點燈用器皿蓋上或放在床底下，乃是放在燈臺上，叫進來的人看見亮光。 17 因為掩藏的事沒有不顯出來的，隱瞞的事沒有不露出來被人知道的。 18 所以，你們應當\textcolor{red}{小心}怎樣聽。因為凡有的，還要加給他；凡沒有的，連他自以為有的也要奪去。


路加福音 21:8-9\\
耶穌說：「你們要\textcolor{red}{謹慎}，不要受迷惑。因為將來有好些人冒我的名來，說『我是基督』，又說『時候近了』。你們不要跟從他們！ 9 你們聽見打仗和擾亂的事，不要驚惶，因為這些事必須先有，只是末期不能立時就到。」



路加福音 21:34-36\\
你们要\textcolor{red}{謹慎}，恐怕因贪食、醉酒并今生的思虑累住你们的心，那日子就如同网罗忽然临到你们， 35 因为那日子要这样临到全地上一切居住的人。 36 你们要时时警醒，常常祈求，使你们能逃避这一切要来的事，得以站立在人子面前。

使徒行傳 13:38-41\\
所以弟兄們，你們當曉得：赦罪的道是由這人傳給你們的！ 39 你們靠摩西的律法，在一切不得稱義的事上信靠這人，就都得稱義了。 40 所以你們務要\textcolor{red}{小心}，免得先知書上所說的臨到你們。 41 主說：『你們這輕慢的人，要觀看，要驚奇，要滅亡！因為在你們的時候，我行一件事，雖有人告訴你們，你們總是不信。』


使徒行傳 20:25-31\\
我素常在你們中間來往，傳講神國的道，如今我曉得，你們以後都不得再見我的面了。 26 所以我今日向你們證明，你們中間無論何人死亡，罪不在我身上， 27 因為神的旨意我並沒有一樣避諱不傳給你們的。 28 聖靈立你們做全群的監督，你們就當為自己\textcolor{red}{謹慎}，也為全群\textcolor{red}{謹慎}，牧養神的教會，就是他用自己血所買來的 d。 29 我知道我去之後，必有凶暴的豺狼進入你們中間，不愛惜羊群。 30 就是你們中間也必有人起來，說悖謬的話，要引誘門徒跟從他們。 31 所以你們應當警醒，記念我三年之久晝夜不住地流淚勸誡你們各人。


林前 3:10-15\\
我照神所給我的恩，好像一個聰明的工頭立好了根基，有別人在上面建造，只是各人要\textcolor{red}{謹慎}怎樣在上面建造。 11 因為那已經立好的根基就是耶穌基督，此外沒有人能立別的根基。 12 若有人用金、銀、寶石、草、木、禾秸在這根基上建造， 13 各人的工程必然顯露；因為那日子要將它表明出來，有火發現，這火要試驗各人的工程怎樣。 14 人在那根基上所建造的工程若存得住，他就要得賞賜； 15 人的工程若被燒了，他就要受虧損，自己卻要得救，雖然得救，乃像從火裡經過的一樣。


林前 10:11-12\\
他們遭遇這些事都要作為鑒戒，並且寫在經上正是警戒我們這末世的人。 12 所以，自己以為站得穩的須要\textcolor{red}{謹慎}，免得跌倒。

林前 16:10-11\\
若是提摩太來到，你們要\textcolor{red}{留心}，叫他在你們那裡無所懼怕，因為他勞力做主的工，像我一樣。 11 所以，無論誰都不可藐視他，只要送他平安前行，叫他到我這裡來，因我指望他和弟兄們同來。

加拉太書 5:13-15
弟兄們，你們蒙召是要得自由，只是不可將你們的自由當做放縱情慾的機會，總要用愛心互相服侍。 14 因為全律法都包在「愛人如己」這一句話之內了。 15 你們要\textcolor{red}{謹慎}，若相咬相吞，只怕要彼此消滅了。

腓立比書 3:2-3
應當\textcolor{red}{防備}犬類，\textcolor{red}{防備}作惡的，\textcolor{red}{防備}妄自行割的。 3 因為真受割禮的，乃是我們這以神的靈敬拜，在基督耶穌裡誇口，不靠著肉體的。


歌羅西書 2:8
你們要\textcolor{red}{謹慎}，恐怕有人用他的理學和虛空的妄言，不照著基督，乃照人間的遺傳和世上的小學，就把你們擄去。


希伯來書 3:12-13
弟兄們，你們要\textcolor{red}{謹慎}，免得你們中間或有人存著不信的惡心，把永生神離棄了。 13 總要趁著還有今日，天天彼此相勸，免得你們中間有人被罪迷惑，心裡就剛硬了。

希伯來書 12:22-25
你們乃是來到錫安山，永生神的城邑，就是天上的耶路撒冷；那裡有千萬的天使， 23 有名錄在天上諸長子之會所共聚的總會，有審判眾人的神和被成全之義人的靈魂， 24 並新約的中保耶穌以及所灑的血。這血所說的比亞伯的血所說的更美。 25 你們總要\textcolor{red}{謹慎}，不可棄絕那向你們說話的。因為那些棄絕在地上警戒他們的，尚且不能逃罪，何況我們違背那從天上警戒我們的呢？



彼得前書 4:7\\
萬物的結局近了，所以你們要你们要\textcolor{red}{謹慎}自守，警醒禱告。

彼得前書 5:8
務要你们要\textcolor{red}{謹守}、警醒，因為你們的仇敵魔鬼如同吼叫的獅子，遍地遊行，尋找可吞吃的人。

約翰二書 1:7-8
因為世上有許多迷惑人的出來，他們不認耶穌基督是成了肉身來的。這就是那迷惑人、敵基督的。 8 你們要\textcolor{red}{小心}，不要失去你們所做的工，乃要得著滿足的賞賜。








